\documentclass{article}
\usepackage{amsmath}
\usepackage{amssymb}

\title{Statistics --- Mathematical Foundations}
\author{Kevin kiplangat}
\date{14-June-26 --present}

\begin{document}
\maketitle

\section{Descriptive Statistics}

\subsection{Mean, Median, and Mode}
The mean, median, and mode are measures of central tendency that summarize a set of data points. The mean is the average of the data points, calculated by summing all the values and dividing by the number of values. The median is the middle value when the data points are arranged in order, and the mode is the value that appears most frequently in the dataset.

The \textbf{mean} is calculated as follows:
\begin{equation}
\bar{x} = \frac{1}{n} \sum_{i=1}^{n} x_i
\end{equation}

where $x_i$ represents each data point and $n$ is the total number of data points.

The \textbf{median} is determined by arranging the data points in ascending order and finding the middle value. If there is an even number of data points, the median is the average of the two middle values. The \textbf{mode} is simply the value that occurs most frequently in the dataset.    


\end{document}